\begin{kbvrule}
\noindent\kbvbullet\hspace{0.5em}ವರ್ಗೀಯ ವ್ಯಂಜನಗಳನ್ನು ಉಚ್ಚಾರದ ಆಧಾರದ ಮೇಲೆ ಅಲ್ಪಪ್ರಾಣ, ಮಹಾಪ್ರಾಣ, ಅನುನಾಸಿಕ ಎಂಬುದಾಗಿ ಮೂರು ಬಗೆಯಾಗಿ ವಿಂಗಡಿಸಬಹುದು.
\end{kbvrule}

\vspace{2.5pt plus 3pt minus 0.5pt}\noindent\hspace*{1.2em}\kbvnum{(೧)}\hspace{0.55em}“ಕಚಟ ತಪ” ಮತ್ತು “ಗಜಡ ದಬ” ಇವು ತೇಲಿಸಿ-ಉಚ್ಚರಿಸ ಸುವ ವ್ಯಂಜನಗಳು ; ಇವುಗಳಿಗೆ “ಅಲ್ಪಪ್ರಾಣ” ಗಳೆಂದು ಹೆಸರು.\par

\vspace{2.5pt plus 3pt minus 0.5pt}\noindent\hspace*{1.2em}\kbvnum{(೨)}\hspace{0.55em}“ಖಛಠ ಥಫ” ಮತ್ತು “ಘ ಝ ಢ ಧ ಭ” - ಇವು ಹೆಚ್ಚು ಒತ್ತು ಕೊಟ್ಟು ಉಚ್ಚರಿಸುವ ವ್ಯಂಜನಗಳು ; ಇವುಗಳಿಗೆ “ಮಹಾಪ್ರಾಣ” ಗಳೆಂದು ಹೆಸರು.\par

\vspace{2.5pt plus 3pt minus 0.5pt}\noindent\hspace*{1.2em}\kbvnum{(೩)}\hspace{0.55em}ಙ ಞ ಣ ನ ಮ - ಇವು ಮೂಗಿನ ಸಹಾಯದಿಂದ ಉಚ್ಚರಿಸ ಸುವ ವ್ಯಂಜನಗಳು ; ಇವುಗಳಿಗೆ “ಅನುನಾಸಿಕ” ಗಳೆಂದು ಹೆಸರು.\par

\begin{kbvrule}
\noindent\kbvbullet\hspace{0.5em}ಅವರ್ಗೀಯ ವ್ಯಂಜನಗಳು : ಯ, ರ, ಲ, ವ, ಶ, ಷ, ಸ, ಹ, ಳ - ಇವು ಅವರ್ಗೀಯ ವ್ಯಂಜನಗಳು ; ಯ-ಕಾರದಿಂದ ಳ-ಕಾರದ ವರೆಗಿನ ಈ ವ್ಯಂಜನಗಳ ಒಟ್ಟು ಸಂಖ್ಯೆ ಒಂಭತ್ತು.
\end{kbvrule}

\kbvsechead{ಯೋಗವಾಹಗಳು :}

\indent ಸ್ವರಾಕ್ಷರದ ಬಲಬದಿಯಲ್ಲಿ ಒಂದು ಸೊನ್ನೆ ಬಂದರೆ “ಅನುಸ್ವಾರ” ವೆನಿಸುವುದು ; ಮತ್ತು ಒಂದರ ಮೇಲೊಂದರಂತೆ ಎರಡು ಸೊನ್ನೆಗಳು ಬಂದರೆ “ವಿಸರ್ಗ”ವೆನಿಸುವುದು. ಈ ಎರಡು ಪ್ರಭೇದಗಳನ್ನು ಯೋಗವಾಹಗಳೆನ್ನುವರು.\par

\indent ಅನುಸ್ವಾರ (ಂ) ; ವಿಸರ್ಗ (ಃ)\par

\indent ಉದಾಹರಣೆಗಳು : ಅಂ ; - ಅಃ\par

\indent ಯೋಗವಾಹಗಳನ್ನು ಸ್ವರಗಳ ಸಹಯೋಗದೊಂದಿಗೆ ಮಾತ್ರ ಉಚ್ಚರಿಸ ಬಹುದಾದ ಕಾರಣ ಅವುಗಳಿಗೆ ಆ ಹೆಸರು ಬಂದಿದೆ.\par

\kbvsechead{ಗುಣಿತಾಕ್ಷರಗಳು :}

\indent ಮೂಲ ವ್ಯಂಜನಗಳು ಅರ್ಧಾಕ್ಷರಗಳು. ಇವುಗಳ ಉಚ್ಚಾರ ಕಠಿಣ ಹಾಗೂ ಅಸ್ಪಷ್ಟ. ಈ ಅಪೂರ್ಣ ಅಕ್ಷರಗಳೊಂದಿಗೆ ಸ್ವರಗಳು ಸೇರಿದಾಗ ಸ್ಪಷ್ಟ ವಾಗಿ ಉಚ್ಚರಿಸಬಹುದಾದ ಹೊಸ ಅಕ್ಷರಗಳು ಲಭಿಸುತ್ತವೆ... ಹೀಗೆ ರೂಪ ತಳೆಯುವ ವ್ಯಂಜನಗಳಿಗೆ “ಗುಣಿತಾಕ್ಷರ” ಗಳೆಂದು ಹೆಸರು.\par
